%% sections/macros.tex
%% Macro condivise — usa \providecommand per evitare conflitti con aastex
%% NON inserire \begin o \end qui

%% ---- Integratori e librerie ----------------------------------------
%% NOTA: \aas è già definita da aastex come abbreviazione di giornale.
%%       Usiamo \AAS per il nostro integratore.
\newcommand{\AAS}{\textsc{AAS}}           % il nostro integratore
\newcommand{\astdyn}{\texttt{AstDyn}}
\newcommand{\saba}{\textsc{SABA4}}
\newcommand{\rkf}{\textsc{RKF7(8)}}
\newcommand{\orbitn}{\texttt{orbitN}}
\newcommand{\rebound}{\texttt{REBOUND}}

%% ---- Unità fisiche --------------------------------------------------
%% \au e \auday sono già in aastex — usiamo versioni estese
%%\newcommand{\unitAU}{\,\mathrm{AU}}
%%newcommand{\unitAUd}{\,\mathrm{AU\,day^{-1}}}
%%\newcommand{\unitAUd2}{\,\mathrm{AU\,day^{-2}}}
\newcommand{\pAU}{\ensuremath{\,\mathrm{AU}}}
\newcommand{\pAUd}{\ensuremath{\,\mathrm{AU\,day^{-1}}}}
\newcommand{\pAUdd}{\ensuremath{\,\mathrm{AU\,day^{-2}}}}
\providecommand{\unitAU}{\pAU}
\providecommand{\unitAUd}{\pAUd}

%% ---- Simboli matematici --------------------------------------------
\newcommand{\dEE}{|\Delta E/E|}
\newcommand{\tauD}{\hat{\tau}_D}
\newcommand{\eps}{\varepsilon}
\newcommand{\rhoH}{\rho(\mathbf{H})}
\newcommand{\Hess}{\mathbf{H}}
\newcommand{\STM}{\boldsymbol{\Phi}}
%% \q e \p collidono con alcuni package — usiamo \vq e \vp (vettori)
\newcommand{\vq}{\mathbf{q}}
\newcommand{\vp}{\mathbf{p}}
\DeclareMathOperator{\tr}{tr}
\newcommand{\Poisson}[2]{\left\{#1,#2\right\}}
\newtheorem{theorem}{Theorem}[section]
\newtheorem{lemma}[theorem]{Lemma}

%% ---- Giornali non in aastex ----------------------------------------
\newcommand{\cmda}{Celest.\ Mech.\ Dyn.\ Astron.}
\newcommand{\physleta}{Phys.\ Lett.\ A}
\newcommand{\cacm}{Commun.\ ACM}
